\section{Related Work}
\label{sec:related}

\paragraph{Biological sequence benchmarks and foundation models.}
DNA benchmarks cover histone marks, promoters, splice sites and regulatory annotations
(GUE/DNABERT-2, 28 tasks \citep{zhou2024dnabert}; Genomic Benchmarks \citep{gresova2022genomic};
DeepSEA \citep{zhou2015deepsea}; DeepSTARR \citep{dealmeida2022deepstarr}); protein
benchmarks cover fluorescence, stability, homology and mutation-fitness assays
(TAPE \citep{rao2019tape}; ProteinGym's 217 DMS assays \citep{notin2023proteingym};
FLIP \citep{dallago2021flip}); localization and function benchmarks
are multi-label (DeepLoc 2.0 \citep{thumuluri2022deeploc}; DeepGOZero
\citep{kulmanov2022deepgozero}); RNA benchmarks cover RBP binding preference
(RNAcompete \citep{ray2013rnacompete}). Foundation models --- ESM-style protein LMs
\citep{lin2023esm2}, protein language tooling \citep{xu2023protst,xu2023biotranslator}, and
multimodal DNA agents \citep{almeida2025chatnt} --- are trained and ranked per family. These
resources are family-siloed: incompatible record formats, split protocols and metric
conventions make cross-family, cross-interface measurement impossible, and most are
single-sequence, single-task-type. BioDecisionBench is not another leaderboard substrate: it
converts these sources into one typed-decision record format with leakage auditing as a
first-class citizen (\cref{sec:admission}) and serves as the evaluation base for the
architecture comparison.

\paragraph{Typed and candidate-scoring decision interfaces.}
The Jev line of systems \citep{almeida2026jev} exposes structured Noul/Choice/Score
probabilistic judgements behind a service interface (weights not public). Open variants
include GLiClass, a ModernBERT uni-encoder with dynamic label sets
\citep{gliclass2025}; NanoJev, a small shared-scorer reimplementation
\citep{nanojev2026}; and SemIf, which reads option logits from a frozen LLM
\citep{semif2025}. Constrained decoding is the generative counterpart
\citep{scholak2021picard}. Each of these validates the interface's feasibility; none
compares \emph{shared scorer vs.\ per-task heads} under matched compute, same data, and
multiple seeds. We report the honest, counter-intuitive outcome of exactly that comparison.

\paragraph{Evaluation methodology and variance.}
Repeated runs and variance reporting are standard advice in ML evaluation
\citep{bouthillier2021variance,miller2024errorbars}; LLM-evaluation work documents
sensitivity to phrasing, prompt format and option order \citep{miller2024errorbars}. In
biology, prior work measured candidate-order sensitivity in the same decision setting:
reordering candidates changes ${\approx}9\%$ of predictions \citep{wang2026biopaws}. This paper adds a concrete \emph{conclusion-reversal}
case for architecture comparison: two runs of one configuration give score Spearman 0.107
vs 0.015 for the shared scorer, flipping the winner; three-seed averaging overturns the
single-run ``shared $\ge$ heads'' verdict (\cref{sec:mainresults}). To our knowledge this is
the first fully documented instance of run variance reversing a decision-architecture
conclusion, and we operationalize the remedy: $\ge3$ seeds, per-task seed means, SEM and
run-std reported separately, and dev early stopping with $\ge256$ dev samples.

\paragraph{Probing and adaptation cost.}
Linear probing and few-shot head fitting are standard tools for characterizing
representations. We operationalize the ``zero-adaptation interface'' claim as a
\emph{catch-up cost}: freeze the encoder, fit a new head for a held-out task with
$N\in\{32,128,512\}$ labels, and report the smallest $N$ at which the head reaches the
shared scorer's zero-shot level (\cref{sec:value}). This turns ``interface flexibility''
from a qualitative claim into a comparable number.
